\clearpage\fancyhead[L]{\small COMET.}\hypertarget{native-338-lower}{}\label{native-338-lower}
\begin{center}\textbf{COMET.}\par\vspace{4pt}\hrule\vspace{4pt}\fontsize{10}{12}\selectfont Astronomisches Jahrbuch for 1814. pp. 254\,..\,256. Berlin 1811.\par\vspace{4pt}\hrule\vspace{4pt}\end{center}

\newcommand{\ObservedNumeratorMark}{\hbox{\vrule width0pt height2.291667bp depth0pt\special{pdf:bcontent}\special{pdf:code q 0 g 1.208333 0.833333 m 1.083333 0.833333 l 1.041667 1.000000 l 0.708333 1.000000 l 0.708333 0.875000 l 0.541667 0.833333 l 0.541667 0.708333 l 0.375000 0.666667 l 0.333333 0.500000 l 0.000000 0.500000 l 0.000000 0.000000 l 0.500000 0.000000 l 0.541667 0.166667 l 0.875000 0.166667 l 0.875000 0.291667 l 1.041667 0.333333 l 1.041667 0.500000 l 1.208333 0.541667 l h 0.000000 2.083333 m 0.000000 1.958333 l 0.125000 1.958333 l 0.166667 1.791667 l 1.041667 1.791667 l 1.083333 1.583333 l 1.375000 1.583333 l 1.375000 2.083333 l 1.208333 2.125000 l 1.208333 2.291667 l 0.166667 2.291667 l 0.166667 2.125000 l h f* Q}\special{pdf:econtent}\kern1.375000bp}}

\begin{flushright}Göttingen, 10 September 1811.\end{flushright}
\par\indent We first saw the comet here at dusk on 22 August, but the view from our observatory is restricted by the city, so we could only make observations starting in September. I used angular distances from fixed stars, which can be observed very well with the sextant mounted on a stand. It has not yet been possible to reduce all the observations made so far. Today I can give only the following three, which are based on distances from $\alpha$ Lyrae and $\alpha$ Aurigae.
\Needspace{9\baselineskip}
\begin{center}\fontsize{10}{12}\selectfont\setlength{\tabcolsep}{4pt}\renewcommand{\arraystretch}{1.15}\begin{tabular}{r|c|c|c}
1811 & Mean time & Right ascension & Northern declination\\
\hline
Sept. 4 & $8^{\mathrm u}\,28^{\mathrm m}\,47^{\mathrm s}$ & $158^\circ\,25'\,24''$ & $39^\circ\,18'\,2''$\\
6 & $8\;48\;38$ & $160\;23\;16$ & $40\;14\;16$\\
7 & $8\;57\;6$ & $161\;26\;10$ & $40\;41\;54$\\
\end{tabular}\end{center}

\par\indent From these observations and new observations by Mr. von \textsc{Zach}, I have improved my provisional parabolic elements, printed in issue 130 of our learned Notices [p. 327 of this volume], as follows.
\Needspace{11\baselineskip}
\begin{center}\fontsize{10}{12}\selectfont\setlength{\tabcolsep}{4pt}\renewcommand{\arraystretch}{1.15}\begin{tabular}{l}
Perihelion passage 1811 Sept. 12\\
\quad $5^{\mathrm u}\,21^{\mathrm m}\,15^{\mathrm s}$ mean time in Göttingen\\
\end{tabular}\end{center}

\begin{center}\fontsize{10}{12}\selectfont\setlength{\tabcolsep}{4pt}\renewcommand{\arraystretch}{1.15}\begin{tabular}{llr}
Logarithm of the least distance & .\;.\;.\;.\;.\;.\;.\;.\;. & $0.017060$\\
Longitude of perihelion & .\;.\;.\;.\;.\;.\;.\;.\;.\;.\;.\;.\;.\;.\;. & $75^\circ\,17'\,34''$\\
Longitude of the ascending node & .\;.\;.\;.\;.\;.\;.\;.\;.\;. & $140\;24\;13$\\
Inclination of the orbit & .\;.\;.\;.\;.\;.\;.\;.\;.\;.\;.\;.\;.\;.\;. & $73\;7\;17$\\
\end{tabular}\end{center}
\noindent Retrograde motion.

\Needspace{19\baselineskip}\SourceLeaf{339}\par\indent Using these elements, I have calculated the geocentric motion for the next 4 months, as shown below.

\begin{center}\fontsize{10}{12}\selectfont\setlength{\tabcolsep}{4pt}\renewcommand{\arraystretch}{1.15}\begin{tabular}{rr|c|c|c|c}
\multicolumn{2}{c|}{Göttingen time} & \shortstack{Right\\ascension} & \shortstack{Northern\\declination} & \shortstack{Log. of distance\\from Earth} & \shortstack{Brightness.}\\
\hline
1811 Sept. 12 & $5^{\mathrm u}$ & $166^\circ\,54'$ & $42^\circ\,58'$ & $0.2079$ & $0.355$\\
17 & $5$ & $173\;46$ & $45\;15$ & $0.1836$ & $0.394$\\
22 & $6$ & $182\;2$ & $47\;17$ & $0.1594$ & $0.432$\\
Oct. 2 & $7$ & $203\;24$ & $49\;34$ & $0.1160$ & $0.490$\\
12 & $8$ & $228\;45$ & $47\;15$ & $0.0911$ & $0.492$\\
22 & $9$ & $251\;23$ & $39\;56$ & $0.0963$ & $0.421$\\
Nov. 1 & $10$ & $267\;58$ & $30\;32$ & $0.1312$ & $0.310$\\
11 & $11$ & $279\;38$ & $22\;18$ & $0.1851$ & $0.208$\\
21 & $12$ & $288\;13$ & $14\;55$ & $0.2415$ & $0.138$\\
Dec. 1 & $13$ & $294\;54$ & $9\;42$ & $0.2976$ & $0.093$\\
11 & $14$ & $300\;24$ & $5\;49$ & $0.3487$ & $0.064$\\
21 & $15$ & $303\;22$ & $3\;0$ & $0.3848$ & $0.048$\\
\end{tabular}\end{center}

\par\indent Because of very unfavorable weather, we saw little of the lunar eclipse on 2 September here, but the occultation of $\lambda$ \emph{Aquarii} that took place during the eclipse (given about 1 hour too early in the Yearbook) was observed very well. Both immersion and emersion took place at the eclipsed but still visible limb of the Moon. The latter could be observed all the better because one of my students had calculated the occultation in advance, so we were prepared for it.

\par\indent Immersion of $\lambda$ Aquarii, mean time $10^{\mathrm u}\,16^{\mathrm m}\,17^{\mathrm s}4$. Emersion $10^{\mathrm u}\,51^{\mathrm m}\,40^{\mathrm s}5$.

\par\indent As you requested, I also enclose a short supplement to my Theoria motus corporum coelestium.
\Needspace{16\baselineskip}\begin{center}\emph{Supplement to Articles 90 and 100 of the Theoria motus corporum coelestium, by the author.}\end{center}
\par\indent In solving the important problem of determining the elliptic or hyperbolic elements from two radius vectors and the angle between them, I have found it very helpful to use an auxiliary quantity, $\xi$ for the ellipse and $\zeta$ for the hyperbola. I appended a table for it to that work. The table was calculated from a continued fraction given there. Its full derivation is not given there, and I have yet to find an opportunity to develop the theory behind it, which is connected with other investigations. Some readers may therefore be glad to find here another way in which that auxiliary quantity could have been calculated just as conveniently.

\par\indent We have (Art. 90)
\[
\xi=x-\frac{5}{6}+\frac{10}{9X}=\frac{xX-\dfrac{5}{6}X+\dfrac{10}{9}}{X}
\]
\par\indent Substituting the series given there for $X$ readily transforms the numerator of this fraction into

\[
\begin{aligned}
&\frac{8}{105}xx\bigl(1+\frac{2\mathord{.}8}{9}x+\frac{3\mathord{.}8\mathord{.}10}{9\mathord{.}11}xx\\
&\qquad+\frac{4\mathord{.}8\mathord{.}10\mathord{.}12}{9\mathord{.}11\mathord{.}13}x^3+\frac{5\mathord{.}8\mathord{.}10\mathord{.}12\mathord{.}14}{9\mathord{.}11\mathord{.}13\mathord{.}15}x^4+\text{etc.}\bigr)
\end{aligned}
\]
\par\indent Thus if we write the series as
\[
1+\frac{2\mathord{.}8}{9}x+\frac{3\mathord{.}8\mathord{.}10}{9\mathord{.}11}xx+\text{etc.}=A
\]
\SourceLeaf{340}\noindent we get
\[
xX-\frac{5}{6}X+\frac{10}{9}=\frac{8}{105}Axx
\]
\[
X=\frac{\dfrac{4}{3}\bigl(1-\dfrac{12}{175}Axx\bigr)}{1-\dfrac{6}{5}x}
\]
\[
\xi=\frac{\dfrac{\mathord{\raisebox{7bp}{\ObservedNumeratorMark}\,}2}{35}Axx\bigl(1-\dfrac{6}{5}x\bigr)}{1-\dfrac{12}{175}Axx}
\]
\noindent This formula always allows us to calculate $\xi$ conveniently and reliably. For $\zeta$ we need only substitute $-z$ for $x$.

\Needspace{8\baselineskip}\par\indent I would only add that $A$ can be calculated even more conveniently by the following formula

\[
\begin{aligned}
A={}&(1-x)^{-\frac{3}{2}}\bigl(1+\frac{1\mathord{.}5}{2\mathord{.}9}x+\frac{1\mathord{.}3\mathord{.}5\mathord{.}7}{2\mathord{.}4\mathord{.}9\mathord{.}11}xx\\
&\qquad+\frac{1\mathord{.}3\mathord{.}5\mathord{.}5\mathord{.}7\mathord{.}9}{2\mathord{.}4\mathord{.}6\mathord{.}9\mathord{.}11\mathord{.}13}x^3+\text{etc.}\bigr)
\end{aligned}
\]
\noindent The derivation of this series from the preceding one, however, rests on considerations that cannot be set out here.
